\section{Sequences as functions}

Sequences are often introduced as lists of numbers:
\[
1,\; \frac12,\; \frac13,\; \frac14,\; \ldots
\]
or
\[
2,\;4,\;6,\;8,\;\ldots
\]
This way of writing sequences is useful, but it can hide the most important idea. A sequence is not a completely different kind of object from a function. A sequence is a special kind of function.

The difference is in the domain.

\begin{definitionbox}
A real sequence is a function
\[
a:\mathbb{N}\to\mathbb{R}.
\]
For each natural number \(n\), the function assigns a real number \(a(n)\). Instead of writing \(a(n)\), we usually write
\[
a_n.
\]
Thus
\[
a_n=a(n).
\]
\end{definitionbox}

The notation changes, but the idea is the same: one input gives one output. For sequences, the inputs are natural numbers. The values are called terms of the sequence.

\subsection*{Terms of a sequence}

If \(a:\mathbb{N}\to\mathbb{R}\), then
\[
a_1,
\quad
 a_2,
\quad
 a_3,
\quad \ldots
\]
are the values of the function at the inputs \(1,2,3,\ldots\). Depending on convention, some sequences start at \(n=0\). What matters is that the starting index must be clear.

\begin{examplebox}
Let
\[
a_n=2n+1,
\qquad n\in\mathbb{N}.
\]
Then
\[
a_1=2\cdot1+1=3,
\]
\[
a_2=2\cdot2+1=5,
\]
and
\[
a_3=2\cdot3+1=7.
\]
The sequence begins
\[
3,\;5,\;7,\;9,\;\ldots
\]
\end{examplebox}

The formula \(a_n=2n+1\) tells us the value of the function \(a\) at the natural number \(n\). It is function notation written in sequence language.

\subsection*{Index and value}

In a sequence, the input is called the index. The output is called the term. These two roles should not be confused.

In the sequence
\[
a_n=2n+1,
\]
the number \(n\) is the index. The number \(a_n\) is the term at index \(n\). For example, when \(n=4\),
\[
a_4=9.
\]
The index is \(4\), and the term is \(9\).

\begin{remarkbox}
The subscript in \(a_n\) is not decoration. It records the input of the sequence-function. The expression \(a_5\) means the value of the sequence at index \(5\).
\end{remarkbox}

\subsection*{Graphs of sequences}

Since a sequence is a function, it can have a graph. But its graph looks different from the graph of a function defined on an interval. For a sequence \(a:\mathbb{N}\to\mathbb{R}\), the graph consists of points
\[
(n,a_n),
\qquad n\in\mathbb{N}.
\]
These points are separated because the inputs are separated.

\begin{examplebox}
For the sequence
\[
a_n=\frac1n,
\qquad n\ge1,
\]
the graph consists of the points
\[
\left(1,1\right),
\quad
\left(2,\frac12\right),
\quad
\left(3,\frac13\right),
\quad
\left(4,\frac14\right),
\quad\ldots
\]
We do not automatically draw a continuous curve through these points, because the sequence is not defined for every real number between \(1\) and \(2\). It is defined only at natural-number inputs.
\end{examplebox}

This is one of the main differences between ordinary real functions on intervals and sequences. A sequence has a discrete domain.

\subsection*{Sequences as ordered data}

Because the domain is \(\mathbb{N}\), the values of a sequence come with a natural order:
\[
a_1,
\quad
a_2,
\quad
a_3,
\quad\ldots
\]
We can ask how the terms change as \(n\) increases. Do they grow? Do they decrease? Do they approach a number? Do they repeat a pattern?

These questions are function questions, but adapted to the special domain \(\mathbb{N}\).

\begin{summarybox}
When reading a sequence, ask:
\begin{itemize}
\item What is the starting index?
\item What is the domain of the sequence?
\item What is the formula or rule for \(a_n\)?
\item What are the first few terms?
\item What is the difference between the index \(n\) and the value \(a_n\)?
\item Should the graph be read as separated points rather than a continuous curve?
\end{itemize}
\end{summarybox}

A sequence is a function with discrete input. This simple fact explains both why sequences belong in a chapter on functions and why they are studied with some different questions.